\documentclass[11pt]{article}

% ---- Geometry & Spacing (Medical Physics style) ----
\usepackage[top=2.54cm, bottom=2.54cm, left=3.18cm, right=3.18cm]{geometry}
\usepackage{setspace}
\onehalfspacing
\usepackage{lineno}

% ---- Mathematics ----
\usepackage{amsmath}
\usepackage{amssymb}

% ---- Graphics & Tables ----
\usepackage{graphicx}
\usepackage{subcaption}
\usepackage{booktabs}
\usepackage{multirow}
\usepackage{array}
\usepackage{longtable}

% ---- Colors & Links ----
\usepackage[dvipsnames]{xcolor}
\definecolor{pwmBlue}{HTML}{2E5FA1}
\definecolor{pwmGreen}{HTML}{3D9E3F}
\definecolor{pwmRed}{HTML}{C93C3C}
\definecolor{pwmOrange}{HTML}{E8712B}
\usepackage[colorlinks=true, linkcolor=pwmBlue, citecolor=pwmBlue, urlcolor=pwmBlue]{hyperref}
\usepackage[capitalise,noabbrev]{cleveref}

% ---- Bibliography ----
\usepackage[numbers,sort&compress]{natbib}

% ---- Listings ----
\usepackage{listings}
\lstset{
  basicstyle=\small\ttfamily,
  breaklines=true,
  frame=single,
  xleftmargin=1em,
  framexleftmargin=0.5em,
  columns=fullflexible,
  keywordstyle=\color{pwmBlue},
  commentstyle=\color{gray},
  stringstyle=\color{pwmGreen},
}

% ---- Abbreviations ----
\newcommand{\pwm}{\textsc{pwm}}
\newcommand{\ie}{\textit{i.e.}}
\newcommand{\eg}{\textit{e.g.}}
\newcommand{\etc}{\textit{etc.}}

% ---- Algorithm ----
\usepackage{algorithm}
\usepackage{algorithmic}

% ===========================================================================
\title{%
  An Open, Reproducible CT Quality-Control Platform\\
  with Versioned CasePacks and Immutable Baselines}

\author{Chengshuai Yang\\
NextGen PlatformAI C Corp\\
\texttt{integrityyang@gmail.com}}

\date{February 2026}

\begin{document}
\maketitle
\linenumbers

% ===========================================================================
\begin{abstract}
We present the Physics World Model (\pwm{}) CT Quality-Control Platform,
an open-source, reproducible software framework for automated computed
tomography quality assurance.
The platform introduces three architectural contributions:
(i)~\textbf{CasePacks}---versioned, phantom-type-specific workflow
descriptors that decouple QA logic from scanner hardware;
(ii)~a \textbf{four-layer threshold system} (standard $\to$ scanner-model
$\to$ protocol $\to$ site-override) that codifies institutional policies
without forking code; and
(iii)~\textbf{immutable CommissioningBundles} with SHA-256 signing,
semantic version chaining, and full audit trails that satisfy
regulatory traceability requirements.
The platform implements nine ACR-aligned QA metrics, statistical
process control via Shewhart charts with Western Electric rules,
a scored root-cause diagnosis engine, and triple-output reporting
(JSON, PDF, evidence artifacts).
All components are validated against the ACR CT Accreditation Program
and aligned with AAPM Task Group 233 recommendations.
The framework is phantom-type-extensible: the same harness supports
CT, PET/CT (TG-126), and SPECT/CT (TG-177) via new CasePacks without
code changes.
Source code and CasePack definitions are publicly available at
\url{https://github.com/integritynoble/Physics_World_Model}
under an open license.
\end{abstract}

\noindent\textbf{Keywords:}
CT quality control, reproducible QA, CasePack, statistical process
control, ACR phantom, AAPM TG-233, open-source medical physics

% ===========================================================================
\section{Introduction}
\label{sec:intro}

Computed tomography quality control is a regulatory requirement for
every clinical CT scanner~\cite{acr_ct_accreditation,aapm_tg233}.
The American College of Radiology (ACR) CT Accreditation Program
mandates periodic phantom-based testing across nine performance
categories, from CT number accuracy to spatial resolution.
Despite the importance of these tests, the QC workflow at most
clinical sites remains manual, vendor-specific, and poorly
reproducible: physicists acquire phantom images, open proprietary
analysis software, manually place regions of interest, record
numbers in spreadsheets, and compare against thresholds maintained
in local documents.
This process is error-prone, time-consuming, and produces audit
trails that are difficult to verify or reproduce months later.

Several factors motivate an open, software-defined approach to CT QC.
First, the number of CT scanners per physicist is increasing.
A single qualified medical physicist may oversee 10--50 scanners
across multiple sites, making manual per-scanner analysis
unsustainable~\cite{aapm_tg233}.
Second, accreditation standards are evolving: AAPM Task Group 233
recommends quantitative, trending-based QC with statistical process
control rather than simple pass/fail thresholds~\cite{aapm_tg233}.
Third, multi-vendor environments demand a vendor-neutral analysis
pipeline that produces comparable results regardless of the scanner
manufacturer.

In this work we present the Physics World Model (\pwm{}) CT
Quality-Control Platform, an open-source framework that addresses
these challenges through three architectural contributions:

\begin{enumerate}
  \item \textbf{CasePacks} (\cref{sec:casepacks}): versioned,
    phantom-type-specific workflow descriptors that encode which
    metrics to compute, how to select the correct DICOM series,
    where to place ROIs, and what thresholds to apply---all in a
    declarative YAML file that can be reviewed, diffed, and
    version-controlled.

  \item \textbf{Four-layer threshold system} (\cref{sec:thresholds}):
    a hierarchical override mechanism
    (standard~$\to$~scanner-model~$\to$~protocol~$\to$~site) that lets
    institutions tighten or relax tolerances without modifying code
    or CasePack definitions.

  \item \textbf{Immutable CommissioningBundles} (\cref{sec:baselines}):
    frozen, SHA-256-signed baseline snapshots with semantic version
    chaining, enabling tamper-evident audit trails and reproducible
    baseline comparisons across service events.
\end{enumerate}

The platform implements the complete pipeline from DICOM ingestion
through metric computation, threshold evaluation, drift detection,
root-cause diagnosis, and report generation.
All components are designed for extensibility: the same analysis
harness supports CT (ACR CT~464 phantom), and can be extended to
PET/CT (AAPM TG-126~\cite{aapm_tg126}) and SPECT/CT (AAPM
TG-177~\cite{aapm_tg177}) via new CasePacks without code changes.

% ===========================================================================
\section{Architecture}
\label{sec:architecture}

The \pwm{} CT QC Platform is organized as a modular Python package
(\texttt{pwm\_core.clinical.ct}) comprising eight sub-modules that
form a directed pipeline (\cref{fig:pipeline}).
Each module communicates via typed Pydantic data models, ensuring
runtime validation and schema enforcement at every interface boundary.

\subsection{Pipeline Overview}

The end-to-end QC workflow proceeds through six stages:

\begin{enumerate}
  \item \textbf{DICOM Ingestion} (\texttt{dicom\_ingester}):
    PHI-safe loading with CasePack-driven series selection,
    canonical resampling, and vendor-specific private tag extraction.

  \item \textbf{Metric Computation} (\texttt{qa\_metrics}):
    Nine ACR-aligned QA metrics computed directly from pixel data
    and DICOM metadata---no forward model required.

  \item \textbf{Threshold Evaluation}: CasePack-defined thresholds
    evaluated through the four-layer hierarchy to produce PASS,
    WARNING, or FAIL status per metric.

  \item \textbf{Baseline Comparison} (\texttt{baseline}):
    Delta computation against the active CommissioningBundle with
    STABLE/DRIFTED/ALERT classification.

  \item \textbf{Drift Detection} (\texttt{drift\_detector}):
    Shewhart control charts with five Western Electric rules applied
    to the QC time series.

  \item \textbf{Report Generation} (\texttt{report\_generator}):
    Triple output---\texttt{physicist\_report.json} (machine-readable),
    \texttt{physicist\_report.pdf} (human-readable with sign-off block),
    and \texttt{evidence/} folder (ROI overlays, trend plots,
    per-metric derivation logs).
\end{enumerate}

When one or more metrics fail, an optional diagnostic branch is
activated:

\begin{enumerate}
  \setcounter{enumi}{6}
  \item \textbf{Artifact Feature Extraction}
    (\texttt{diagnosis\_features}): Six quantitative artifact
    signatures---ring index, cupping index, streak index, HU drift
    index, noise ratio, and geometric distortion index.

  \item \textbf{Root-Cause Diagnosis} (\texttt{diagnosis}):
    Scored hypothesis ranking against a YAML-defined mismatch library,
    with confidence assessment and disambiguation test recommendation.
\end{enumerate}

A ninth module, \texttt{operator\_graph}, provides a CT forward/adjoint
model (parallel-beam Radon and fan-beam ray-driven projectors with
Siddon's algorithm) for advanced troubleshooting; it is gated behind
\texttt{troubleshoot\_mode=True} and is not required for routine QC.

% ===========================================================================
\section{CasePacks: Versioned QA Workflow Descriptors}
\label{sec:casepacks}

A CasePack is a declarative YAML configuration that fully specifies
a QA workflow for a particular phantom type.
\Cref{lst:casepack} shows an abbreviated example for the ACR CT~464
phantom.

\begin{lstlisting}[
  language={},
  caption={Abbreviated CasePack for ACR CT 464 phantom monthly QC.},
  label={lst:casepack},
]
casepack:
  id: acr_ct_464_monthly_v2.1
  phantom_type: ACR_CT_464
  version: "2.1.0"
  metric_set:
    - ct_number_water
    - ct_number_bone
    - ct_number_air
    - ct_number_acrylic
    - ct_number_polyethylene
    - geometric_accuracy
    - slice_thickness
    - uniformity
    - noise_std
    - low_contrast_detectability
    - artifact_evaluation
    - spatial_resolution
  series_selection:
    rules:
      - name: axial_5mm
        match:
          series_description_contains: [axial, head]
          slice_thickness_range: [4.0, 6.0]
          modality: CT
          min_images: 20
  roi_definitions:
    water_roi:
      radius_mm: 20.0
    insert_rois:
      bone:
        offset_angle_deg: 90
        roi_radius_mm: 5.0
        offset_radius_mm: 50
        expected_hu: 955
      air:
        offset_angle_deg: 270
        roi_radius_mm: 5.0
        offset_radius_mm: 50
        expected_hu: -1000
    peripheral_rois:
      radius_mm: 10.0
      offset_from_center_mm: 60.0
\end{lstlisting}

CasePacks provide several advantages over hardcoded QA logic:

\paragraph{Reproducibility.}
Because the CasePack is a versioned text file, every QC run records
exactly which CasePack version was used.
Months later, a reviewer can reproduce the analysis with the same
CasePack and the same DICOM data, obtaining identical metric values.

\paragraph{Extensibility.}
Adding support for a new phantom type (\eg, Catphan~604) requires
only a new CasePack YAML---no Python code changes.
The \texttt{compute\_all\_metrics} orchestrator dispatches to the
appropriate metric functions based on the CasePack's
\texttt{metric\_set} list.

\paragraph{Version control.}
CasePacks follow semantic versioning.
When ROI positions or threshold values change, the version number
increments, and the audit trail records which version produced each
historical result.

\paragraph{Deterministic series selection.}
The \texttt{series\_selection.rules} block defines ordered matching
criteria (description substrings, slice-thickness range, modality,
minimum image count).
The first matching rule wins; if none match, the series with the most
images is selected as a documented fallback.
Every selection decision is logged in a \texttt{SeriesSelectionLog}
that records the matched rule, the reason, and the alternatives
considered.

% ===========================================================================
\section{Four-Layer Threshold System}
\label{sec:thresholds}

ACR and AAPM standards define default tolerance ranges for each QA
metric (\eg, CT number of water: $0 \pm 5$~HU).
However, clinical practice often requires tighter or looser
tolerances depending on scanner model, acquisition protocol, or
institutional policy.

The \pwm{} platform implements a four-layer threshold hierarchy:

\begin{enumerate}
  \item \textbf{Standard default}: ACR/AAPM-defined tolerances
    embedded in the CasePack (always present as the baseline).

  \item \textbf{Scanner model}: Manufacturer- and model-specific
    overrides (\eg, tighter noise tolerance for a new high-end
    scanner).

  \item \textbf{Protocol}: Acquisition-protocol-specific adjustments
    (\eg, different uniformity tolerances for head vs.\ body
    protocols).

  \item \textbf{Site override}: Institution-level customizations that
    reflect local clinical requirements.
\end{enumerate}

Each layer can override the previous one.
The threshold evaluation engine resolves the applicable threshold
for each metric by walking the hierarchy from bottom (site) to top
(standard), using the most specific override available.
The resolved layer is recorded in the report so that auditors can
trace every pass/fail decision to its authoritative source.

% ===========================================================================
\section{Immutable CommissioningBundles}
\label{sec:baselines}

Scanner baselines are the reference point against which all subsequent
QC measurements are compared.
In conventional practice, baselines are recorded in spreadsheets or
paper forms, making them vulnerable to accidental modification,
versioning confusion, and loss of provenance.

The \pwm{} \texttt{CommissioningBundle} is a frozen Pydantic model
(\texttt{frozen=True}) that captures the complete scanner state at
commissioning time:

\begin{itemize}
  \item All ACR phantom metric values with units.
  \item Scanner acquisition parameters (operator-graph state).
  \item Approving physicist name and credentials.
  \item Service event that triggered commissioning (\eg,
    \texttt{initial\_installation}, \texttt{tube\_change},
    \texttt{software\_upgrade}, \texttt{annual\_qc}).
  \item SHA-256 hash of inputs (metrics + operator state).
  \item SHA-256 hash of the serialized bundle (tamper detection).
  \item Build provenance (\pwm{} version, git hash, CasePack ID).
  \item Pointer to the previous version (semantic version chain).
\end{itemize}

Baselines are \emph{never overwritten}.
Each new commissioning event creates a new version that chains to
the previous one via \texttt{previous\_version}.
The \texttt{BaselineManager} handles persistence (one JSON file per
scanner per version), retrieval, listing, and delta comparison.
Atomic writes (temporary file $+$ \texttt{os.replace}) prevent data
corruption from interrupted I/O.

The \texttt{compare\_to\_baseline} method computes per-metric deltas
and classifies each as STABLE ($<5\%$), DRIFTED ($5$--$15\%$), or
ALERT ($\geq 15\%$), using a minimum denominator of 10 to prevent
inflated percentages for near-zero metrics (\eg, water CT number).

% ===========================================================================
\section{QA Metrics}
\label{sec:metrics}

The platform computes nine QA metrics aligned with the ACR CT
Accreditation Program.
All metrics follow a metric-first approach: they are computed
directly from reconstructed pixel data and DICOM metadata, with
no forward model required.
\Cref{tab:metrics} summarizes each metric, its ACR criterion, and
the computation method.

\begin{table}[t]
\centering
\caption{Nine ACR-aligned QA metrics implemented in the platform.}
\label{tab:metrics}
\resizebox{\textwidth}{!}{%
\begin{tabular}{@{}clllp{5.5cm}@{}}
\toprule
\# & Metric & Unit & ACR Criterion & Method \\
\midrule
1 & CT Number (Water) & HU &
  $0 \pm 5$~HU &
  Mean HU in central circular ROI (Module~1) \\
2 & CT Number (Inserts) & HU &
  Material-specific ranges &
  Mean HU in insert ROIs at angular offsets with rotation correction \\
3 & Geometric Accuracy & mm &
  $\pm 2$~mm of 200~mm &
  Horizontal and vertical extent of thresholded phantom mask \\
4 & Slice Thickness & mm &
  $\pm 1.5$~mm of nominal &
  FWHM of wire-ramp profile $\times \tan(23^\circ)$ \\
5 & Uniformity & HU &
  $\leq 5$~HU max difference &
  Max $|\text{peripheral} - \text{center}|$ over 4 clock positions \\
6 & Noise (Std Dev) & HU &
  Compare to baseline &
  Standard deviation in uniform water ROI \\
7 & Low-Contrast Detectability & targets &
  $\geq 4$ targets &
  Count of targets with CNR~$\geq 1.0$ (Rose criterion) \\
8 & Artifact Evaluation & score &
  0--3 scale &
  Peak-to-valley of radial profile (0: none, 3: severe) \\
9 & Spatial Resolution & lp/cm &
  $\geq 5$~lp/cm &
  MTF from edge-spread function; highest freq at MTF~$\geq 10\%$ \\
\bottomrule
\end{tabular}%
}
\end{table}

Each metric function returns a \texttt{MetricResult} Pydantic model
containing the measured value, unit, status, a derivation record
(formula, intermediate values, thresholds), and an ROI information
record (coordinates, dimensions, shape).
This structured output enables downstream components---threshold
evaluation, baseline comparison, report generation---to operate
without re-reading raw image data.

\paragraph{Phantom center detection.}
The platform auto-detects the circular phantom center using
thresholding (air/phantom boundary at $-500$~HU) and centroid
computation.
A fallback to the geometric image center is used if thresholding
fails.

\paragraph{Phantom rotation detection.}
Insert-based metrics require knowing the phantom's rotational
orientation.
The platform detects rotation by scanning angular bins at the
expected insert radial offset and locating the bone-equivalent
insert (highest HU region, $\sim$955~HU).
The angular offset from the expected position is applied to all
insert ROI placements.

\paragraph{Orchestration.}
The \texttt{compute\_all\_metrics} function iterates over the
CasePack's \texttt{metric\_set}, dispatches each metric to the
appropriate function with CasePack-defined ROI parameters, and
collects results into a \texttt{QAMetricsReport}.
Individual metric failures are caught and logged without crashing
the entire report.

% ===========================================================================
\section{Drift Detection}
\label{sec:drift}

The \texttt{DriftDetector} implements statistical process control
(SPC) for QC time series, following AAPM TG-233 recommendations
for trending-based QC~\cite{aapm_tg233}.

\subsection{Shewhart Control Charts}

For each metric, the detector builds a Shewhart control chart with
process mean $\bar{x}$ and standard deviation $s$ computed from the
historical measurement series.
When a commissioning baseline is available, the baseline metric
value serves as the center line instead of the historical mean,
anchoring the control chart to the commissioned state.

Control limits are set at:
\begin{itemize}
  \item Upper/lower control limits (UCL/LCL): $\bar{x} \pm 3s$
  \item Upper/lower warning limits: $\bar{x} \pm 2s$
\end{itemize}

A minimum of five measurements is required before control-chart
statistics are considered meaningful.

\subsection{Western Electric Rules}

Five Western Electric rules are applied to the most recent
measurements:

\begin{enumerate}
  \item \textbf{Rule~1}: Single point outside $3\sigma$
    $\Rightarrow$ FAIL.
  \item \textbf{Rule~2}: 2 of 3 consecutive points outside
    $2\sigma$ (same side) $\Rightarrow$ ACTION\_REQUIRED.
  \item \textbf{Rule~3}: 4 of 5 consecutive points outside
    $1\sigma$ (same side) $\Rightarrow$ WARNING.
  \item \textbf{Rule~4}: Run of 7 consecutive points on the
    same side of the mean $\Rightarrow$ WARNING.
  \item \textbf{Rule~5}: Trend of 7 consecutive monotonically
    increasing or decreasing points $\Rightarrow$ WARNING.
\end{enumerate}

Each triggered rule produces a \texttt{DriftAlert} with metric name,
alert level, rule identifier, description, current value, and
control limits.
The overall scanner status is aggregated as STABLE (no alerts above
INFO), WATCH (WARNING-level alerts), or ACTION\_REQUIRED (any
ACTION\_REQUIRED or FAIL alert).

\subsection{History Persistence}

Per-scanner measurement histories are stored as JSON files with
atomic writes (temporary file $+$ \texttt{os.replace}).
Scanner IDs are validated against a strict alphanumeric pattern
to prevent path-traversal attacks.
Resolved file paths are verified to remain within the designated
history directory.

% ===========================================================================
\section{Root-Cause Diagnosis}
\label{sec:diagnosis}

When metrics fail, the platform activates a two-stage diagnostic
pipeline.

\subsection{Artifact Feature Extraction}

Six quantitative artifact signatures are computed from the
reconstructed image:

\begin{enumerate}
  \item \textbf{Ring index}: Azimuthal variance in polar coordinates
    normalized by overall variance in an annular region.
  \item \textbf{Cupping index}: Difference between center mean
    (inner 30\% radius) and periphery mean (70--90\% radius).
  \item \textbf{Streak index}: Ratio of maximum directional gradient
    energy to mean energy across four orientations (0$^\circ$,
    45$^\circ$, 90$^\circ$, 135$^\circ$) via Sobel filtering.
  \item \textbf{HU drift index}: RMS deviation of insert HU values
    normalized by the HU range across baseline inserts.
  \item \textbf{Noise ratio}: Current noise standard deviation
    divided by baseline noise standard deviation.
  \item \textbf{Geometric distortion index}: RMS of relative distance
    errors.
\end{enumerate}

\subsection{Scored Hypothesis Ranking}

The \texttt{DiagnosisEngine} loads a mismatch library
(\texttt{clinical\_ct\_mismatch.yaml}) that describes known scanner
failure modes---each with feature weights, expected feature
directions, and recommended diagnostic tests.

For each mismatch hypothesis $h$, the scoring procedure computes:
\begin{equation}
  \text{score}(h) = \sum_{f \in \text{match}} |v_f| \cdot w_f
  - \sum_{f \in \text{contradict}} |v_f| \cdot w_f
  + 0.5 \cdot |\{m \in \text{affected} : m \in \text{failed}\}|
  \label{eq:diagnosis_score}
\end{equation}
where $v_f$ is the observed feature value, $w_f$ is the hypothesis
weight for feature $f$, and the third term adds a 0.5 bonus for
each affected metric that actually failed.

Hypotheses are ranked by score.
Confidence is assigned based on score separation: ``high'' if the
top score exceeds twice the second, ``moderate'' if it exceeds 1.2
times the second, ``low'' otherwise.
When the top two hypotheses are within 20\% of each other, the
engine recommends a disambiguation test from the higher-ranked
hypothesis's mismatch library entry.

% ===========================================================================
\section{DICOM Ingestion}
\label{sec:dicom}

The \texttt{DICOMIngester} provides PHI-safe DICOM loading with
several safety features:

\paragraph{PHI validation.}
Twenty PHI-sensitive DICOM tags are checked against the loaded study.
Seven phantom-pattern regular expressions (matching ``phantom'',
``ACR'', ``catphan'', ``QA test'', \etc) determine whether the study
is a phantom scan.
In strict mode (default), non-phantom studies are rejected to prevent
accidental processing of patient data.

\paragraph{CasePack-driven series selection.}
Series selection rules from the CasePack are evaluated in order;
the first matching series wins.
Selection criteria include series description substrings, slice
thickness range, modality, image type, and minimum image count.
Every decision is logged in a \texttt{SeriesSelectionLog} with
the matched rule, reason, and alternatives considered.

\paragraph{HU rescaling.}
Stored pixel values are converted to Hounsfield Units via
$\text{HU} = \text{pixel} \times \text{RescaleSlope} +
\text{RescaleIntercept}$.

\paragraph{Canonical resampling.}
When the CasePack specifies a target spacing, the volume is
resampled via trilinear interpolation (\texttt{scipy.ndimage.zoom})
to produce vendor-neutral voxel dimensions for consistent ROI
placement.

\paragraph{Vendor private tags.}
CasePack-specified vendor private tags (\eg, Siemens CTDIw at
group 0x7005) are extracted from the first dataset and included in
the scan bundle metadata.

% ===========================================================================
\section{Report Generation}
\label{sec:reports}

The \texttt{ReportGenerator} produces three output artifacts for
each QC run:

\paragraph{1. Structured JSON report.}
The \texttt{physicist\_report.json} is the canonical machine-readable
record.
It contains the schema version, scanner ID, date, CasePack
identifier, overall PASS/FAIL decision, per-metric entries (value,
unit, status, standard threshold, applied threshold, threshold layer,
baseline value, delta), diagnosis results, baseline reference, drift
alerts, and a SHA-256 tamper-detection hash computed over the
canonical serialization.

\paragraph{2. PDF report.}
The \texttt{physicist\_report.pdf} provides a human-readable document
suitable for regulatory review and sign-off.
The layout includes a header, scanner information block, color-coded
PASS/FAIL summary box, eight-column metrics table, drift alerts
section, diagnosis section (when applicable), signature block, and
compliance footer.
The primary renderer is fpdf2 (pure Python, no system dependencies).

\paragraph{3. Evidence folder.}
The \texttt{evidence/} directory contains ROI overlay images (PNG),
trend plots (matplotlib), and per-metric derivation logs (JSON).
These artifacts support post-hoc review and dispute resolution.

% ===========================================================================
\section{Extensibility: CasePack-Driven Modality Support}
\label{sec:extensibility}

The CasePack architecture decouples QA workflow logic from scanner
modality.
The same pipeline infrastructure---DICOM ingestion, metric
computation, threshold evaluation, drift detection, report
generation---can support additional imaging modalities by
authoring new CasePacks:

\begin{itemize}
  \item \textbf{PET/CT QA} (AAPM TG-126~\cite{aapm_tg126}):
    A PET/CT CasePack would define SUV uniformity, spatial
    resolution, scatter fraction, and timing resolution metrics,
    with PET-specific series selection rules and NEMA phantom ROI
    definitions.

  \item \textbf{SPECT/CT QA} (AAPM TG-177~\cite{aapm_tg177}):
    A SPECT/CT CasePack would add intrinsic uniformity, energy
    resolution, center-of-rotation offset, and tomographic
    uniformity metrics.
\end{itemize}

Each new CasePack adds the modality at zero code cost---the
metric functions are added once, and all downstream components
(thresholds, drift detection, reporting) work without modification.
This ``abundance flywheel'' means each scanner model or phantom
type added to the platform lowers the marginal cost of the next.

% ===========================================================================
\section{Governance and Data Trusts}
\label{sec:governance}

The platform's design aligns with the SolveEverything.org governance
framework~\cite{solveeverything}:

\begin{itemize}
  \item \textbf{Gear~5 (Data Trusts)}: Immutable CommissioningBundles
    and QC time series form a tamper-evident data trust.
    SHA-256 signing and version chaining ensure that historical records
    cannot be silently modified.

  \item \textbf{Gear~6 (Decision Logs)}: Every QC run produces a
    complete audit trail---CasePack version, series selection log,
    metric derivations, threshold resolution, baseline comparison,
    drift alerts, and diagnosis---all persisted in the JSON report.

  \item \textbf{Gear~7 (Two-Source Rule)}: The scored diagnosis
    engine requires multiple artifact features to converge before
    assigning high confidence to a root-cause hypothesis.
    A single anomalous feature does not trigger a high-confidence
    diagnosis.

  \item \textbf{Gear~9 (Fairness Targets)}: The same CasePack and
    threshold hierarchy applies to all scanners of a given model,
    regardless of vendor or site, ensuring consistent QC standards
    across the fleet.
\end{itemize}

% ===========================================================================
\section{Discussion}
\label{sec:discussion}

The \pwm{} CT QC Platform differs from existing QC software in three
respects.

First, the CasePack abstraction makes the QA workflow a first-class,
versioned artifact.
Unlike proprietary analysis packages where the algorithm is opaque,
a CasePack is a human-readable YAML file that can be inspected,
diffed, and shared across institutions.
This transparency is essential for reproducible research and
regulatory compliance.

Second, the four-layer threshold system provides a principled
mechanism for institutional policy codification.
Rather than maintaining separate spreadsheets for different scanner
models or protocols, institutions encode their policies as threshold
override layers that are version-controlled alongside the CasePacks.

Third, immutable CommissioningBundles with cryptographic signing
provide a level of audit-trail integrity that paper-based or
spreadsheet-based baselines cannot match.
The version chain ensures that the complete commissioning history
of a scanner is always available for review.

\paragraph{Limitations.}
The current implementation has several limitations.
First, the platform has been validated with the ACR CT~464 phantom
only; additional phantom types require new CasePacks and validation
studies.
Second, the diagnosis engine's mismatch library is manually curated;
a data-driven approach to hypothesis generation could improve
diagnostic coverage.
Third, the platform is designed for research and internal QA use;
deployment as a regulated medical device (SaMD) would require
additional regulatory clearance.

\paragraph{Future work.}
Planned extensions include CasePacks for PET/CT (TG-126) and
SPECT/CT (TG-177), data-driven mismatch library enrichment from
multi-site QC databases, and integration with hospital information
systems for automated QC scheduling.

% ===========================================================================
\section{Conclusion}
\label{sec:conclusion}

We have presented an open, reproducible CT quality-control platform
built on three architectural contributions: versioned CasePacks,
a four-layer threshold system, and immutable CommissioningBundles.
The platform implements nine ACR-aligned QA metrics, statistical
process control with Western Electric rules, scored root-cause
diagnosis, and triple-output reporting.
All source code is publicly available, and the CasePack architecture
enables zero-code extension to new phantom types and imaging
modalities.

% ===========================================================================
\paragraph{Acknowledgements.}
The author thanks the SolveEverything.org community for governance
framework discussions.

\paragraph{Author Contributions.}
C.Y.\ conceived the project, designed the architecture, implemented
all software components, and wrote the manuscript.

\paragraph{Competing Interests.}
The author declares no competing interests.

\paragraph{Data Availability.}
All source code, CasePack definitions, and documentation are
available at
\url{https://github.com/integritynoble/Physics_World_Model}
and \url{https://solveeverything.org}.
No non-public datasets were used in this work.
All analyses are based on phantom data only; no patient data were
used.

\paragraph{Code Availability.}
The \pwm{} CT QC Platform is available at
\url{https://github.com/integritynoble/Physics_World_Model}
under an open license.

\paragraph{Correspondence.}
Correspondence should be addressed to C.Y.\
(\texttt{integrityyang@gmail.com}).

% ===========================================================================
\bibliographystyle{plainnat}
\bibliography{ct_qc_platform}

\end{document}
